\SourceLeaf{152}
\[
\begin{array}{@{}l@{\;}l@{}}
\text{summa duorum terminorum}&=\displaystyle\frac{n-1}{h}-\frac{(n-h-1)(n-h)}{hn}\\[7pt]
\text{summa trium terminorum}&=\displaystyle\frac{n-1}{h}-\frac{(n-h-1)(n-h)(n-h+1)}{hn(n+1)}\quad\text{etc.}
\end{array}
\]
et quum pars altera (per II) formet seriem ultra omnes limites decrescentem,
summa illa in infinitum continuata statui debet $=\frac{n-1}{h}$. Hinc $N+N'+N''$ etc.
in infinitum continuata semper manebit minor quam $\frac{N(n-1)}{h}$, et proin $M+M'+M''$
etc. certo ad summam finitam converget. Q.~E.~D.

VI. Ut ea, quae generaliter de serie $M,M',M''$ etc. demonstravimus, ad
coëfficientes potestatum $x^m,x^{m+1},x^{m+2}$ etc. in serie $F(\alpha,\beta,\gamma,x)$, applicentur,
statuere oportebit $\lambda=2$, $A=\alpha+\beta$, $B=\alpha\beta$, $a=\gamma+1$, $b=\gamma$, unde quinque assertiones in art. praec. propositae sponte demanant.

\SixthArticle{17}
Disquisitio itaque de indole summae seriei $F(\alpha,\beta,\gamma,1)$ natura sua restringitur ad casum, quo $\gamma-\alpha-\beta$ est quantitas positiva, ubi illa summa semper
exhibet quantitatem finitam. Praemittimus autem observationem sequentem. Si
coëfficientes seriei $1+ax+bxx+cx^3+$ etc. $=S$ inde a certo termino ultra
omnes limites decrescunt, productum
\[
(1-x)S=1+(a-1)x+(b-a)xx+(c-b)x^3+\quad\text{etc.}
\]
pro $x=1$ statuere oportet $=0$, etiamsi summa ipsius seriei $S$ infinite magna
evadat. Quoniam enim collectis duobus terminis summa fit $=a$, collectis tribus $=b$, collectis quatuor $=c$ etc., limes summae in infinitum continuatae
est $=0$. Quoties itaque $\gamma-\alpha-\beta$ est quantitas positiva, statuere oportet
$(1-x)F(\alpha,\beta,\gamma-1,x)=0$ pro $x=1$, unde per aequationem 15 art. 7 habemus
\[
0=\gamma(\alpha+\beta-\gamma)F(\alpha,\beta,\gamma,1)+(\gamma-\alpha)(\gamma-\beta)F(\alpha,\beta,\gamma+1,1),\quad\text{sive}
\]
\[
[37]\qquad F(\alpha,\beta,\gamma,1)=\frac{(\gamma-\alpha)(\gamma-\beta)}{\gamma(\gamma-\alpha-\beta)}F(\alpha,\beta,\gamma+1,1)
\]
Quare quum perinde fiat
\[
\begin{aligned}
F(\alpha,\beta,\gamma+1,1)&=\frac{(\gamma+1-\alpha)(\gamma+1-\beta)}{(\gamma+1)(\gamma+1-\alpha-\beta)}F(\alpha,\beta,\gamma+2,1)\\[6pt]
F(\alpha,\beta,\gamma+2,1)&=\frac{(\gamma+2-\alpha)(\gamma+2-\beta)}{(\gamma+2)(\gamma+2-\alpha-\beta)}F(\alpha,\beta,\gamma+3,1)
\end{aligned}
\]
et sic porro, erit generaliter, $k$ denotante integrum positivum quemcunque
\SourceLeaf{153}
\[
\begin{array}{@{}ll@{}}
&F(\alpha.\beta,\gamma,1)\quad\text{aequalis producto ex}\quad F(\alpha,\beta,\gamma+k,1)\\[6pt]
\text{in}&(\gamma-\alpha)(\gamma+1-\alpha)(\gamma+2-\alpha)\mathinner{.\,.\,.\,.}(\gamma+k-1-\alpha)\\[5pt]
\text{in}&(\gamma-\beta)(\gamma+1-\beta)(\gamma+2-\beta)\mathinner{.\,.\,.\,.}(\gamma+k-1-\beta)\\[8pt]
\multicolumn{2}{@{}l@{}}{\text{diviso per productum}}\\[6pt]
\text{ex}&\gamma(\gamma+1)(\gamma+2)\mathinner{.\,.\,.\,.}(\gamma+k-1)\\[5pt]
\text{in}&(\gamma-\alpha-\beta)(\gamma+1-\alpha-\beta)(\gamma+2-\alpha-\beta)\mathinner{.\,.\,.}(\gamma+k-1-\alpha-\beta)
\end{array}
\]
\SixthArticle{18}
Introducamus abhinc sequentem notationem:
\[
[38]\qquad \Pi(k,z)=\frac{1\mathbin{.}2\mathbin{.}3\mathinner{.\,.\,.\,.\,.}k}{(z+1)(z+2)(z+3)\mathinner{.\,.\,.\,.}(z+k)}k^z
\]
ubi $k$ natura sua subintelligitur designare integrum positivum, qua restrictione
$\Pi(k,z)$ exhibet functionem duarum quantitatum $k,z$ prorsus determinatam. Hoc
modo facile intelligitur, theorema in fine art. praec. propositum ita exhiberi posse
\[
[39]\qquad F(\alpha,\beta,\gamma,1)=\frac{\Pi(k,\gamma-1)\mathbin{.}\Pi(k,\gamma-\alpha-\beta-1)}{\Pi(k,\gamma-\alpha-1)\mathbin{.}\Pi(k,\gamma-\beta-1)}\mathbin{.}F(\alpha,\beta,\gamma+k,1)
\]
\SixthArticle{19}
Operae pretium erit, indolem functionis $\Pi(k,z)$ accuratius perpendere.
Quoties $z$ est integer negativus, functio manifesto valorem infinite magnum obtinet, simulac ipsi $k$ tribuitur valor satis magnus. Pro valoribus integris
ipsius $z$ non negativis autem habemus
\[
\begin{aligned}
\Pi(k,0)&=1\\[5pt]
\Pi(k,1)&=\frac{1}{1+\frac1k}\\[5pt]
\Pi(k,2)&=\frac{1\mathbin{.}2}{(1+\frac1k)(1+\frac2k)}\\[5pt]
\Pi(k,3)&=\frac{1\mathbin{.}2\mathbin{.}3}{(1+\frac1k)(1+\frac2k)(1+\frac3k)}
\end{aligned}
\]
etc. sive generaliter
\[
[40]\qquad \Pi(k,z)=\frac{1\mathbin{.}2\mathbin{.}3\mathinner{.\,.\,.\,.\,.\,.}z}{(1+\frac1k)(1+\frac2k)(1+\frac3k)\mathinner{.\,.\,.\,.}(1+\frac zk)}
\]
\SourceLeaf{154}
Generaliter autem pro \textit{quovis} valore ipsius $z$ habemus
\[
[41]\qquad \Pi(k,z+1)=\Pi(k,z)\mathbin{.}\frac{1+z}{1+\frac{1+z}{k}}
\]
\[
[42]\qquad \Pi(k+1,z)=\Pi(k,z)\mathbin{.}\left\{\frac{(1+\frac1k)^{z+1}}{1+\frac{1+z}{k}}\right\}
\]
adeoque, quum $\Pi(1,z)=\frac1{z+1}$,
\[
[43]\qquad \Pi(k,z)=\frac1{z+1}\mathbin{.}\frac{2^{z+1}}{1^z\mathbin{.}(2+z)}\mathbin{.}\frac{3^{z+1}}{2^z(3+z)}\mathbin{.}\frac{4^{z+1}}{3^z(4+z)}\mathinner{.\,.\,.\,.}\frac{k^{z+1}}{(k-1)^z(k+z)}
\]
\SixthArticle{20}
Imprimis vero attentione dignus est \textit{limes}, ad quem pro valore dato ipsius
$z$ functio $\Pi(k,z)$ continuo converget, dum $k$ in infinitum crescit. Sit primo $h$
valor finitus ipsius $k$ maior quam $z$, patebitque, si $k$ transire supponatur ex $h$ in
$h+1$, logarithmum ipsius $\Pi(k,z)$ accipere incrementum, quod per seriem convergentem sequentem exprimitur
\[
\frac{z(1+z)}{2(h+1)^2}+\frac{z(1-zz)}{3(h+1)^3}+\frac{z(1+z^3)}{4(h+1)^4}+\frac{z(1-z^4)}{5(h+1)^5}+\quad\text{etc.}
\]
Si itaque $k$ e valore $h$ transit in $h+n$, logarithmus ipsius $\Pi(k,z)$ accipiet incrementum
\[
\begin{aligned}
&\tfrac12 z(1+z)\left(\frac1{(h+1)^2}+\frac1{(h+2)^2}+\frac1{(h+3)^2}+\text{etc.}+\frac1{(h+n)^2}\right)\\[5pt]
&+\tfrac13 z(1-zz)\left(\frac1{(h+1)^3}+\frac1{(h+2)^3}+\frac1{(h+3)^3}+\text{etc.}+\frac1{(h+n)^3}\right)\\[5pt]
&+\tfrac14 z(1+z^3)\left(\frac1{(h+1)^4}+\frac1{(h+2)^4}+\frac1{(h+3)^4}+\text{etc.}+\frac1{(h+n)^4}\right)\\[5pt]
&+\text{etc.}
\end{aligned}
\]
quod semper finitum manere, etiamsi $n$ in infinitum crescat, facile demonstrari
potest. Quare nisi iam in $\Pi(h,z)$ factor infinitus affuerit, i. e nisi $z$ sit numerus integer negativus, limes ipsius $\Pi(k,z)$ pro $k=\infty$ certo erit quantitas finita. Manifesto itaque $\Pi(\infty,z)$ tantummodo a $z$ pendet, sive functionem ipsius
$z$ ex asse determinatam exhibet, quae abhinc simpliciter per $\Pi z$ denotabitur.
Definimus itaque functionem $\Pi z$ per valorem producti
\[
\frac{1\mathbin{.}2\mathbin{.}3\mathinner{.\,.\,.\,.\,.}k\mathbin{.}k^z}{(z+1)(z+2)(z+3)\mathinner{.\,.\,.\,.}(z+k)}
\]
\SourceLeaf{155}
pro $k=\infty$, aut si mavis per limitem producti infiniti
\[
\frac{1}{z+1}.\frac{2^{z+1}}{1^z(2+z)}.
\frac{3^{z+1}}{2^z(3+z)}.\frac{4^{z+1}}{3^z(4+z)}\quad\mathrm{etc.}
\]

\SixthArticle{21}
Ex aequatione 41 statim sequitur aequatio fundamentalis
\[
\begin{array}{@{}ll@{}}
[44]&\displaystyle\Pi(z+1)=(z+1)\Pi z
\end{array}
\]
unde generaliter, designante $n$ integrum positivum quemcunque
\[
\begin{array}{@{}ll@{}}
[45]&\displaystyle\Pi(z+n)=(z+1)(z+2)(z+3)\,.\,.\,.\,.(z+n)\Pi z
\end{array}
\]
Pro valore integro negativo ipsius $z$ erit valor functionis $\Pi z$
infinite magnus; pro valoribus integris non negativis habemus
\[
\begin{array}{rcl}
\Pi0&=&1\\
\Pi1&=&1\\
\Pi2&=&2\\
\Pi3&=&6\\
\Pi4&=&24\quad\mathrm{etc.}
\end{array}
\]
atque generaliter
\[
\begin{array}{@{}ll@{}}
[46]&\displaystyle\Pi z=1.2.3\,.\,.\,.\,.z
\end{array}
\]
Sed male haec proprietas functionis nostrae tamquam ipsius definitio
venditatur, quippe quae natura sua ad valores integros restringitur,
et praeter functionem nostram infinitis aliis
(e.\,g.\ $\cos2\pi z.\Pi z$, $\cos\pi z^{2n}\Pi z$ etc., denotante $\pi$
semiperipheriam circuli, cuius radius $=1$) communis est.

\SixthArticle{22}
Functio $\Pi(k,z)$, etiamsi generalior videatur quam $\Pi z$, tamen
abhinc nobis superflua erit, quum facile ad posteriorem reducatur.
Colligitur enim e combinatione aequationum 38, 45, 46
\[
\begin{array}{@{}ll@{}}
[47]&\displaystyle\Pi(k,z)=\frac{k^z\Pi k.\Pi z}{\Pi(k+z)}
\end{array}
\]
Ceterum nexus harum functionum cum iis, quas clar.\ \textsc{Kramp}
\emph{facultates nu\SourceLeaf{156}mericas} nominavit, per se obvius est.
Scilicet facultas numerica, quam hic auctor per
$a^{bIc}$
designat, in signis nostris est


\[
=\frac{c^b b^{\frac{a}{c}-1}\Pi b}{\Pi(b,\frac{a}{c}-1)}
=\frac{c^b\Pi(\frac{a}{c}+b-1)}{\Pi(\frac{a}{c}-1)}
\]
Sed consultius videtur, functionem \emph{unius} variabilis in analysin
introducere, quam functionem trium variabilium, praesertim quum hanc
ad illam reducere liceat.

\SixthArticle{23}
Continuitas functionis $\Pi z$ interrumpitur, quoties ipsius valor fit
infinite magnus, i.\,e. pro valoribus integris negativis ipsius $z$.
Erit itaque illa positiva a $z=-1$ usque ad $z=\infty$, et quum pro
utroque limite $\Pi z$ obtineat valorem infinite magnum, inter ipsos
dabitur valor minimus, quem esse $=0{,}8856024$ atque respondere valori
$z=0{,}4616321$ invenimus. Inter limites $z=-1$ et $z=-2$, valor
functionis $\Pi z$ fit negativus, inter $z=-2$ atque $z=-3$ iterum
positivus et sic porro, uti ex aequ.\ 44 sponte sequitur. Porro patet,
si omnes valores functionis $\Pi z$ inter limites arbitrarios unitate
differentes e.\,g. a $z=0$ usque ad $z=1$ pro notis habere liceat,
valorem functionis pro quovis alio valore reali ipsius $z$ adiumento
aequationis 45 facile inde deduci posse. Ad hunc finem construximus
\emph{tabulam}, ad calcem huius sectionis annexam, quae ad figuras
viginti exhibet logarithmos briggicos functionis $\Pi z$, pro $z=0$
usque ad $z=1$ per singulas partes centesimas summa cura computatos,
ubi tamen monendum, figuram ultimam vigesimam interdum una duabusve
unitatibus erroneam esse posse.

\SixthArticle{24}
Quum limes functionis $F(\alpha,\beta,\gamma+k,1)$, crescente $k$ in
infinitum, manifesto sit unitas, aequatio 39 transit in hanc
\[
\begin{array}{@{}ll@{}}
[48]&\displaystyle F(\alpha,\beta,\gamma,1)=
\frac{\Pi(\gamma-1).\Pi(\gamma-\alpha-\beta-1)}
{\Pi(\gamma-\alpha-1).\Pi(\gamma-\beta-1)}
\end{array}
\]
quae formula exhibet solutionem completam quaestionis, quae obiectum
huius sectionis constituit. Sponte hinc sequuntur aequationes elegantes:
\SourceLeaf{157}
\[
\begin{array}{@{}ll@{}}
[49]&\displaystyle F(\alpha,\beta,\gamma,1)=F(-\alpha,-\beta,\gamma-\alpha-\beta,1)\\[3pt]
[50]&\displaystyle F(\alpha,\beta,\gamma,1).F(-\alpha,\beta,\gamma-\alpha,1)=1\\[3pt]
[51]&\displaystyle F(\alpha,\beta,\gamma,1).F(\alpha,-\beta,\gamma-\beta,1)=1
\end{array}
\]
in quarum prima $\gamma$, in secunda $\gamma-\beta$, in tertia
$\gamma-\alpha$ debet esse quantitas positiva.

\SixthArticle{25}
Applicemus formulam 48 ad quasdam ex aequationibus art.\ 5.
Formula XIII, statuendo $t=90^0=\frac12\pi$, fit
$\frac12\pi=F(\frac12,\frac12,\frac32,1)$, sive aequivalet aequationi notae
\[
\frac12\pi=1+\frac{1.1}{2.3}+\frac{1.1.3}{2.4.5}
+\frac{1.1.3.5}{2.4.6.7}+\mathrm{etc.}
\]
Quare quum per formulam 48 habeatur
$F(\frac12,\frac12,\frac32,1)=\frac{\Pi\frac12.\Pi(-\frac12)}{\Pi0.\Pi0}$,
atque sit $\Pi0=1$, $\Pi\frac12=\frac12\Pi(-\frac12)$, fit
$\pi=(\Pi(-\frac12))^2$ sive
\[
\begin{array}{@{}ll@{}}
[52]&\displaystyle\Pi(-\tfrac12)=\BareRoot{\pi}\\[3pt]
[53]&\displaystyle\Pi\tfrac12=\tfrac12\BareRoot{\pi}
\end{array}
\]
Formula XVI art.\ 5, quae aequivalet aequationi notae
\[
\sin nt=n\sin t-\frac{n(nn-1)}{2.3}\sin t^3
+\frac{n(nn-1)(nn-9)}{2.3.4.5}\sin t^5-\mathrm{etc.}
\]
atque generaliter pro quovis valore ipsius $n$ locum habet, si modo
$t$ limites $-90^0$ et $+90^0$ non transgrediatur, dat pro $t=\frac12\pi$
\[
\sin\frac{n\pi}{2}=
\frac{n\Pi\frac12.\Pi(-\frac12)}{\Pi(-\frac12n).\Pi\frac12n}
\]
unde deducitur formula elegans
\[
\Pi\frac12n.\Pi(-\frac12n)=\frac{\frac12n\pi}{\sin\frac12n\pi},
\quad\hbox{sive statuendo }n=2z
\]
\[
\begin{array}{@{}ll@{}}
[54]&\displaystyle\Pi(-z).\Pi(+z)=\frac{z\pi}{\sin z\pi}\\[3pt]
[55]&\displaystyle\Pi(-z).\Pi(z-1)=\frac{\pi}{\sin z\pi}
\end{array}
\]
nec non scribendo $z+\frac12$ pro $z$
\[
\begin{array}{@{}ll@{}}
[56]&\displaystyle\Pi(-\tfrac12+z).\Pi(-\tfrac12-z)=\frac{\pi}{\cos z\pi}
\end{array}
\]
\SourceLeaf{158}
E combinatione formulae 54 cum definitione functionis $\Pi$ sequitur,
$\frac{z\pi}{\sin z\pi}$ esse limitem producti infiniti
\[
\frac{(1.2.3.4\,.\,.\,.\,.k)^2}
{(1-zz)(4-zz)(9-zz)\,.\,.\,.\,.(kk-zz)}
\]
crescente $k$ in infinitum, adeoque
\[
\sin z\pi=z\pi(1-zz)(1-\frac{zz}{4})(1-\frac{zz}{9})\quad
\mathrm{etc.\ in\ inf.}
\]
similique modo ex 56 deducitur
\[
\cos z\pi=(1-4zz)(1-\frac{4zz}{9})(1-\frac{4zz}{25})\quad
\mathrm{etc.\ in\ inf.}
\]
formulae notissimae, quae ab analystis per methodos prorsus diversas
erui solent.

\SixthArticle{26}
Designante $n$ numerum integrum, valor expressionis
\[
\frac{n^{nz}\Pi(k,z).\Pi(k,z-\frac1n).\Pi(k,z-\frac2n)
\,.\,.\,.\,.\Pi(k,z-\frac{n-1}{n})}{\Pi(nk,nz)}
\]
rite collectus invenitur
\[
=\frac{(1.2.3\,.\,.\,.\,.k)^n n^{nk}}
{1.2.3\,.\,.\,.\,.nk.k^{\frac12(n-1)}}
\]
adeoque a $z$ est independens, sive idem manebit, quicunque valor ipsi
$z$. tribuatur. Exhiberi poterit itaque, quoniam $\Pi(k,0)=\Pi(nk,0)=1$,
per productum
\[
\Pi(k,-\frac1n).\Pi(k,-\frac2n).\Pi(k,-\frac3n)
\,.\,.\,.\,.\Pi(k,-\frac{n-1}{n})
\]
Crescente igitur $k$ in infinitum, nanciscimur
\[
\begin{gathered}
\frac{n^{nz}\Pi z.\Pi(z-\frac1n).\Pi(z-\frac2n)
\,.\,.\,.\,.\Pi(z-\frac{n-1}{n})}{\Pi nz}\\[4pt]
=\Pi(-\frac1n).\Pi(-\frac2n).\Pi(-\frac3n)
\,.\,.\,.\,.\Pi(-\frac{n-1}{n})
\end{gathered}
\]
Productum ad dextram, in se ipsum ordine factorum inverso multiplicatum,
producit, per form.\ 55,
\[
\frac{\pi}{\sin\frac1n\pi}.\frac{\pi}{\sin\frac2n\pi}.
\frac{\pi}{\sin\frac3n\pi}\,.\,.\,.\,.
\frac{\pi}{\sin\frac{n-1}{n}\pi}=\frac{(2\pi)^{n-1}}{n}
\]
Unde habemus theorema elegans
\SourceLeaf{159}
\[
\begin{aligned}
{}\text{[57]}\quad &
\frac{n^{nz}\Pi z.\Pi(z-\frac1n).\Pi(z-\frac2n).\,.\,.\,.
\Pi(z-\frac{n-1}{n})}{\Pi nz}
=\frac{(2\pi)^{\frac12(n-1)}}{\BareRoot n}
\end{aligned}
\]

\SixthArticle{27}
Integrale $\int x^{\lambda-1}(1-x^\mu)^\nu\,\mathrm{d}x$,
ita acceptum, ut evanescat pro $x=0$, exprimitur per seriem
sequentem, siquidem $\lambda,\mu$ sunt quantitates positivae:
\[
\frac{x^\lambda}{\lambda}
-\frac{\nu x^{\mu+\lambda}}{\mu+\lambda}
+\frac{\nu(\nu-1)x^{2\mu+\lambda}}{1.2.(2\mu+\lambda)}
-\mathrm{etc.}
=\frac{x^\lambda}{\lambda}F(-\nu,\frac\lambda\mu,
\frac\lambda\mu+1,x^\mu)
\]
Hinc ipsius valor pro $x=1$ erit
\[
=\frac{\Pi\frac\lambda\mu.\Pi\nu}
{\lambda\Pi(\frac\lambda\mu+\nu)}.
\]
Ex hoc theoremate omnes relationes, quas ill. \textsc{Euler} olim
multo labore evolvit, sponte demanant. Ita e. g. statuendo
\[
\int\frac{\mathrm{d}x}{\BareRoot{(1-x^4)}}=A,\qquad
\int\frac{xx\,\mathrm{d}x}{\BareRoot{(1-x^4)}}=B
\]
erit $A=\frac{\Pi\frac14.\Pi(-\frac12)}{\Pi(-\frac14)}$,
$B=\frac{\Pi\frac34.\Pi(-\frac12)}{3\Pi\frac14}
=\frac{\Pi(-\frac14).\Pi(-\frac12)}{4\Pi\frac14}$,
adeoque $AB=\frac14\pi$. Simul hinc sequitur, quoniam
$\Pi\frac14.\Pi(-\frac14)=\frac{\frac14\pi}{\sin\frac14\pi}
=\frac\pi{\BareRoot{8}}$,
\[
\Pi\frac14=\BareRoot[4]{(\frac18\pi AA)}
=\BareRoot[4]{\frac{\pi^3}{128BB}},\qquad
\Pi(-\frac14)=\BareRoot[4]{\frac{\pi^3}{8AA}}
=\BareRoot[4]{(2\pi BB)}
\]
Valor numericus ipsius $A$, computante \textsc{Stirling}, habetur
$=1{,}3110287771\;\allowbreak 4605987$, valor ipsius $B$, secundum eundem
auctorem, $=0{,}5990701173\;\allowbreak 6779611$, ex nostro calculo,
artificio peculiari innixo, $=0{,}5990701173\;\allowbreak 6779610372$.

Generaliter facile ostendi potest, valorem functionis $\Pi z$, si $z$
sit quantitas rationalis $=\frac m\mu$, denotantibus $m,\mu$
integros, ex $\mu-1$ valoribus determinatis talium integralium pro
$x=1$ deduci posse, et quidem permultis modis diversis. Accipiendo
enim pro $\lambda$ numerum integrum atque pro $\nu$ fractionem,
cuius denominator $=\mu$, valor illius integralis semper reducitur
ad tres $\Pi z$, ubi $z$ est fractio cum denominatore $=\mu$;
quodvis vero huiusmodi $\Pi z$ vel ad $\Pi(-\frac1\mu)$,
vel ad $\Pi(-\frac2\mu)$, vel ad $\Pi(-\frac3\mu)$ etc.
vel ad $\Pi(-\frac{\mu-1}{\mu})$ reduci potest per formulam 45,
siquidem $z$ revera est fractio; si enim $z$ est integer, $\Pi z$
per se constat. Ex illis vero integralium valoribus, generaliter
loquendo, quodvis $\Pi(-\frac m\mu)$,
\SourceLeaf{160}
si $m<\mu$, per eliminationem erui potest\HistoricalNote{17}.
Quin adeo semissis talium integralium sufficiet, si formulam 54
simul in auxilium vocamus. Ita e. g. statuendo
\[
\begin{gathered}
\int\frac{\mathrm{d}x}{\BareRoot[5]{(1-x^5)}}=C,\quad
\int\frac{\mathrm{d}x}{\BareRoot[5]{(1-x^5)^2}}=D,\quad
\\\int\frac{\mathrm{d}x}{\BareRoot[5]{(1-x^5)^3}}=E,\quad
\int\frac{\mathrm{d}x}{\BareRoot[5]{(1-x^5)^4}}=F,
\end{gathered}
\]
erit
\[
\begin{gathered}
C=\Pi\frac15.\Pi(-\frac15),\quad
D=\frac{\Pi\frac15.\Pi(-\frac25)}{\Pi(-\frac15)},\quad
\\E=\frac{\Pi\frac15.\Pi(-\frac35)}{\Pi(-\frac25)},\quad
F=\frac{\Pi\frac15.\Pi(-\frac45)}{\Pi(-\frac35)}
\end{gathered}
\]
Hinc propter $\Pi\frac15=\frac15\Pi(-\frac45)$, habemus
\[
\begin{gathered}
\Pi(-\frac15)=\BareRoot[5]{\frac{5C^4}{DEF}},\quad
\Pi(-\frac25)=\BareRoot[5]{\frac{25C^3D^3}{EEFF}},\quad
\\\Pi(-\frac35)=\BareRoot[5]{\frac{125CCDDEE}{F^3}},\\
\Pi(-\frac45)=\BareRoot[5]{(625CDEF)}
\end{gathered}
\]
Formulae 54, 55 adhuc suppeditant
\[
C=\frac\pi{\sin\frac15\pi},\qquad
\frac DF=\frac{\sin\frac15\pi}{\sin\frac25\pi}
\]
ita ut duo integralia $D,E$, vel $E$ et $F$ sufficiant, ad omnes
valores $\Pi(-\frac15)$, $\Pi(-\frac25)$ etc. computandos.

\SixthArticle{28}
Statuendo $y=\nu x$, atque $\mu=1$,
$\frac{\Pi\lambda.\Pi\nu}{\lambda\Pi(\lambda+\nu)}$
erit valor integralis
$\int\frac{y^{\lambda-1}(1-\frac y\nu)^\nu\,\mathrm{d}y}{\nu^\lambda}$ ab $y=0$ usque ad $y=\nu$, sive valor integralis
$\int y^{\lambda-1}(1-\frac y\nu)^\nu\,\mathrm{d}y$
inter eosdem limites
$=\frac{\nu^\lambda\Pi\lambda.\Pi\nu}{\lambda\Pi(\lambda+\nu)}
=\frac{\Pi(\nu,\lambda)}\lambda$ (form. 47), siquidem $\nu$
denotet integrum. Iam crescente $\nu$ in infinitum, limes ipsius
$\Pi(\nu,\lambda)$ erit $=\Pi\lambda$, limes ipsius
$(1-\frac y\nu)^\nu$ autem $e^{-y}$, denotante $e$ basin
logarithmorum hyperbolicorum. Quamobrem si $\lambda$ est positiva,
$\frac{\Pi\lambda}\lambda$ sive $\Pi(\lambda-1)$ exprimet integrale
$\int y^{\lambda-1}e^{-y}\,\mathrm{d}y$ ab $y=0$ usque ad
$y=\infty$, sive scribendo $\lambda$ pro $\lambda-1$,
$\Pi\lambda$ est valor integralis
$\int y^\lambda e^{-y}\,\mathrm{d}y$ ab $y=0$ usque ad
$y=\infty$, si $\lambda+1$ est quantitas positiva.

Generalius statuendo $y=z^\alpha$,
$\alpha\lambda+\alpha-1=\beta$, transit
$\int y^\lambda e^{-y}\,\mathrm{d}y$ in
$\int\alpha z^\beta e^{-z^\alpha}\,\mathrm{d}z$, quod itaque
inter limites $z=0$ atque $z=\infty$ sumtum exprimetur per
$\Pi(\frac{\beta+1}\alpha-1)$

\noindent sive

Valor integralis $\int z^\beta e^{-z^\alpha}\,\mathrm{d}z$,
a $z=0$ usque ad $z=\infty$ fit
$=\frac{\Pi(\frac{\beta+1}\alpha-1)}\alpha
=\frac{\Pi\frac{\beta+1}\alpha}{\beta+1}$
si modo $\alpha$ atque $\beta+1$ sunt quantitates positivae
(si utraque est negativa, in\SourceLeaf{161}tegrale per $-\frac{\Pi\frac{\beta+1}\alpha}{\beta+1}$
exprimetur). Ita e. g. pro $\beta=0$, $\alpha=2$, valor
integralis $\int e^{-zz}\,\mathrm{d}z$ invenitur
$=\Pi\frac12=\frac12\BareRoot\pi$.

\SixthArticle{29}
Ill. \textsc{Euler} pro summa logarithmorum
$\log1+\log2+\log3+\mathrm{etc.}+\log z$ eruit seriem
$(z+\frac12)\log z-z+\frac12\log2\pi
+\frac{\mathfrak A}{1.2z}-\frac{\mathfrak B}{3.4z^3}
+\frac{\mathfrak C}{5.6z^5}-\mathrm{etc.}$
ubi $\mathfrak A=\frac16$, $\mathfrak B=\frac1{30}$,
$\mathfrak C=\frac1{42}$ etc. sunt numeri \textsc{Bernoulliani}.
Per hanc itaque seriem exprimitur $\log\Pi z$; etiamsi enim primo
aspectu haec conclusio ad valores integros restricta videatur,
tamen rem propius contemplando invenietur, evolutionem ab
\textsc{Eulero} adhibitam (Instit. Calc. Diff. Cap.\,\textsc{vi}. 159)
saltem ad valores positivos fractos eodem iure applicari posse,
quo ad integros: supponit enim tantummodo functionem ipsius $z$,
in seriem evolvendam, esse talem, ut ipsius diminutio, si $z$
transeat in $z-1$, exhiberi possit per theorema \textsc{Taylori},
simulque ut eadem diminutio sit $=\log z$. Conditio prior innititur
\emph{continuitati} functionis, adeoque locum non habet pro
valoribus negativis ipsius $z$, ad quos proïn seriem illam
extendere non licet: conditio posterior autem functioni
$\log\Pi z$ generaliter competit sine restrictione ad valores
integros ipsius $z$. Statuamus itaque
\[
\begin{aligned}
{}\text{[58]}\quad &\log\Pi z=(z+\frac12)\log z-z+\frac12\log2\pi\\[4pt]
&+\frac{\mathfrak A}{1.2z}-\frac{\mathfrak B}{3.4z^3}
+\frac{\mathfrak C}{5.6z^5}-\frac{\mathfrak D}{7.8z^7}
+\mathrm{etc.}
\end{aligned}
\]
Quum hinc quoque habeatur
\[
\begin{aligned}
&\log\Pi2z=(2z+\frac12)\log2z-2z+\frac12\log2\pi\\[4pt]
&+\frac{\mathfrak A}{1.2.2z}-\frac{\mathfrak B}{3.4.8z^3}
+\frac{\mathfrak C}{5.6.32z^5}-\frac{\mathfrak D}{7.8.128z^7}
+\mathrm{etc.}
\end{aligned}
\]
atque per formulam 57, statuendo $n=2$,
\[
\log\Pi(z-\frac12)=\log\Pi2z-\log\Pi z
-(2z+\frac12)\log2+\frac12\log2\pi,\quad\text{fit}
\]
\[
\begin{aligned}
{}\text{[59]}\quad &\log\Pi(z-\frac12)=z\log z-z+\frac12\log2\pi\\[4pt]
&-\frac{\mathfrak A}{1.2.2z}+\frac{7\mathfrak B}{3.4.8z^3}
-\frac{31\mathfrak C}{5.6.32z^5}+\frac{127\mathfrak D}{7.8.128z^7}
-\mathrm{etc.}
\end{aligned}
\]
Hae duae series pro valoribus magnis ipsius $z$ ab initio satis
promte convergunt, ita ut summam approximatam commode satisque
exacte colligere liceat: attamen probe notandum est, pro quovis
valore dato ipsius $z$, quantumvis magno, praecisionem limitatam
tantummodo obtineri posse, quum numeri \textsc{Bernoulliani}
seriem hypergeometricam constituant, adeoque series illae, si modo
satis longe extendantur, certo e convergentibus divergentes evadant.
Ceterum negari nequit, theoriam talium serierum divergentium adhuc
quibusdam difficultatibus premi, de quibus forsan alia occasione
pluribus commentabimur.
